\honest{Superhero Bias}{Eliezer Yudkowsky}{2007}

An armed kidnapper holds hostages and has announced that he will start killing. Real police rarely kick down the door, but sometimes they must. In one world, a superhero who can throw cars and whose skin annihilates bullets storms the room to save two hundred children. In another, a New York police officer storms it to save three prostitutes. Who is the greater hero, and who gets the comic book?

The halo effect is that ``perceptions of all positive traits are correlated.'' \nb{The review behind the usual list of traits found the effect near zero for honesty and kindness; see ``The Halo Effect.''} So a strong, invulnerable hero also seems braver. \nb{No study is cited for this. The support is a line from a comic.} As Adam Warren's \textsc{Empowered} asks, how tough can it be to act brave when you are nearly invulnerable?

Fame ``seems to combine additively'' with every other trait; I cannot remember whether I read this or thought of it. Gandhi was protected by his celebrity. The people marching with him, who could be beaten or shot without headlines, took more risk than he did. \nb{At Dharasana in 1930 marchers were clubbed while Gandhi was in jail. But his prominence also made him a target: a mob attacked him in 1897, and there were several attempts on his life before he was killed in 1948. The post counts only the protection.} This is no criticism of Gandhi, whose actions took courage, though less than marching anonymously. The point of nonviolent resistance is not to show courage; going over Niagara Falls in a barrel would do that more easily. The bias is that his fame gets added to his real virtue. When you think of nonviolence you think of Gandhi, not of a marcher who was beaten, walked with a limp for the rest of her life, and whose name no one remembers.

Is it greater to risk your life for two hundred or for three? As an outcome, two hundred: ``Whoever saves a single life, it is as if he had saved the whole world'' is ``terrible moral advice if you've got to pick one or the other.'' \nb{The saying is from the Mishnah, where it warns witnesses in capital cases; it was not offered as advice about whom to save.} As revealed virtue, whoever would risk their life for three shows more courage than whoever would do it only for two hundred.

But choosing the smaller rescue to show virtue would be ``the moral equivalent of manslaughter.'' Someone who risks their life in order to be virtuous reveals far less than someone who does it to save others. You cannot reveal virtue by trying to. Choosing a dangerous way to save the world over a safe one does not make you a hero, since wanting to look like a hero is a lost purpose. Virtuous people seek safer ways to save more, and so reveal less. That is confusing but not contradictory. Whatever risk remains is real courage. The officer, with no powers and no protection from bullets, ``reveals far greater virtue than Superman.''
